\chapter{Голос и модель}
\label{ch:voice}

\section{Субстрат model-agnostic; личность — трансформ}

Память, настроение и граф — состояния \emph{вне} языковой модели. Значит субстрат
\emph{model-agnostic}: его можно строить на любой базовой модели и заменить базу
позже. Это де-рискует проблему потолка малой модели: когнитивная машинерия переживёт
смену базы.

Если факты и здравый смысл лежат в графе (дёшево на диске/в RAM), базовой модели не
нужно быть энциклопедией. Её роль — \emph{лингвистический процессор и голос}: часть
«Rinness» живёт в ней. Можно взять малую дообученную модель ($3$--$7$B в квантовании)
вместо тяжёлой $8$B-универсальной; факты уходят в граф, ёмкость модели
концентрируется на голосе.

\begin{pitfall}[«Меньше — строго лучше» слишком оптимистично]
Малая модель слабее в \emph{рассуждении} и связности: снимаешь фактическую нагрузку,
но не reasoning/coherence. Наблюдение из практики: рантайм prompt-Рин уже процентов
на шестьдесят состоит из анти-повторной обвязки \emph{на} $8$B — вниз с голосом
хуже. Разумный ориентир — $3$--$7$B, не $1$B, и это надо \emph{померить}, а не
постулировать.
\end{pitfall}

\section{Инструмент $\varphi$: скорер по осям характера}

Центральный инструмент — \emph{скорер характера} $\varphi$ (измерительный прибор:
строишь прибор — и он де-конфузит весь отбор). Берём $d$ осей (из архетипики:
аффективная температура, колкость, теплота, агентность, участие-против-сервиса, тип
близости, тип юмора, стиль конфликта). На каждую — \emph{рубрика} (описание $+$
якорь для оценок $0$ и $10$). Тогда $\varphi(\text{текст})\in\mathbb{R}^d$ — прогон
LLM-судьи. Отдельно задаём целевой $\mathbf{t}$. На старте — LLM-as-judge;
классификаторы по осям позже.

\section{Суперпозиция муз — это отбор в признаковом пространстве}

Каждого персонажа можно оценить по каждой оси — персонаж есть \emph{точка} в
$d$-мерном пространстве характера. Музы — существующие персонажи, свои точки.

\emph{Наивная суперпозиция} (обучение на смеси полных корпусов) даёт среднее в
\emph{центроиде} облака муз, а не в $\mathbf{t}$: $\mathbb{E}[\varphi]=\text{centroid}\ne\mathbf{t}$.
Центроид не есть Рин; интересна муза именно отклонением от центроида (центроид яркой
цундере мало отличается от безымянной героини). Усреднение вдобавок размывает
отличительные черты (большая дисперсия по осям).

\emph{Правильно} — построить выборку с $\mathbb{E}[\varphi]\approx\mathbf{t}$: не
сваливать корпуса, а \emph{отбирать} примеры (реплики, сцены), чтобы их агрегатный
центр сидел в $\mathbf{t}$. Технически — отбор данных / rejection sampling под целевой
стилевой вектор: (1)~операционализируем оси; (2)~скорим $\varphi(\text{пример})$;
(3)~отбираем/перевзвешиваем, пока $\mathbb{E}[\varphi]\approx\mathbf{t}$. Короче: не
«смешать музы» (размытый центроид), а «отфильтровать корпус к $\mathbf{t}$»
(выраженная Рин).

\section{Пайплайн по стадиям}

\begin{center}
\begin{tikzpicture}[font=\scriptsize,>=Stealth,node distance=6mm,
  b/.style={draw,rounded corners,align=center,minimum height=8mm,text width=20mm}]
  \node[b] (s1) {Стадия 1\\ сбор реплик муз};
  \node[b,right=of s1] (s2) {Стадия 2\\ отбор по $\varphi\!\to\!\mathbf{t}$};
  \node[b,right=of s2] (s3) {Стадия 3\\ LoRA: голос};
  \node[b,right=of s3] (s4) {Стадия 4\\ ситуация$\to$реакция};
  \node[b,right=of s4] (s5) {Стадия 5\\ замкнуть на граф};
  \draw[->] (s1)--(s2); \draw[->] (s2)--(s3); \draw[->] (s3)--(s4); \draw[->] (s4)--(s5);
  \draw[->] (s5.south) .. controls +(0,-0.8) and +(0,-0.8) .. (s3.south)
     node[midway,below,font=\scriptsize]{консолидация $\to$ новые данные};
\end{tikzpicture}
\end{center}

\paragraph{Стадия 1 — сбор.} Субтитры (\texttt{.srt},\texttt{.ass}), фан-скрипты,
quote-вики, сценарии VN. Извлекаем строки целевых персонажей. Узкое место —
\emph{атрибуция говорящего}. Идея кормить видео каптионеру возможна, но окупается
позже и лишь ради данных о \emph{физическом} поведении: для голоса субтитры дают
почти всё дешевле.

\paragraph{Стадия 2 — отбор (фильтрацией, не написанием).} $\varphi$ по каждой
реплике; оставляем близкие к $\mathbf{t}$. Это корпус голоса.

\paragraph{Стадия 3 — голос (LoRA).} Дообучаем малую модель на отобранных репликах —
языковое моделирование на тексте-в-стиле. Ситуации \emph{не нужны}. LoRA обучает
низкоранговую добавку $\Delta W=BA$, $\mathrm{rank}(BA)=r\ll\dim$, малую по памяти;
влезает на $8$\,ГБ для малой модели.

\paragraph{Стадия 4 — выбор реакции.} Ситуации нужны, эмбеддинги мостят несовпадение.
\emph{(A)~Синтетический} (реком.): ситуации из базы смысла $\to$ промпт сильной модели
в голосе $\to$ реакции $\to$ фильтр $\varphi$ — пары (ситуация$\to$реакция).
\emph{(B)~Майнинг+retrieval}: из сцен муз пары (ситуация, реакция), эмбеддинг ситуаций,
в рантайме ближайший сосед. Именно здесь «ситуации могут не совпадать, эмбеддинги
справятся».

\paragraph{Стадия 5 — замкнуть на граф.} Малая голос-модель $+$ граф; граф даёт
контекст (mood-смещённый retrieval), модель генерит в голосе; выходы (фильтр $\varphi$
$+$ значимость \eqref{eq:salience}) консолидируются обратно.

\section{Bootstrap: генератор $\ne$ деплой-модель}

Как получить характер-образующий граф без готовой Рин? Финальная Рин не нужна —
зачаток есть: \emph{prompt-Рин} (модульные промпты $+$ базовая модель).

\begin{enumerate}
\item строим нейтральный эмбриональный граф (глава~\ref{ch:experiment});
\item селим prompt-Рин как \emph{bootstrap-генератор}: живёт на графе, генерит
характерные реакции;
\item \emph{консолидация пишет отпечатки обратно} — граф становится character-shaped
через \emph{использование}, а не ручным reorganization-проходом;
\item на накопленном поведении дообучаем малую модель.
\end{enumerate}

\begin{pitfall}[Стабильность и потолок]
Аудит логов: prompt-Рин — \emph{Рин с перебоями}. Ядро иногда пробивает хорошо
(реплики, дословно совпадающие с установкой мотивации), но свободный чат деградирует
в липкие повторы, часть ответов нарушает явные правила (запрещённый «перформанс
трагедии», сбои рода), и \emph{guided} чище \emph{raw} — по внутреннему критерию
«ядро промпта недостаточно сильно» (потолок $8$B). Вывод: как \emph{безфильтровый}
bootstrap она шумна — вобьёшь и тики. Два хода: (i)~сидировать не в real-time,
сильной моделью (CPU-offload/квантизация или разовый API), а малую сажать на
обогащённый граф — генератор и деплой \emph{разные}; (ii)~фильтровать выходы $\varphi$.
\end{pitfall}

\section{Raw против guided и псевдо-DPO}

\emph{Raw} менее стабилен, но на пике \emph{ближе} к Рин; \emph{guided} стабилен, но
«ассистентский» (стабильно мимо цели). Для деплоя важна стабильность, для сидирования
— нет: нужны \emph{сэмплы пиковой Rinness}. Низкий потолок фильтрацией не поднять, а
нестабильность — можно (сэмплируй много, фильтруй $\varphi$). Для bootstrap: raw $+$
oversample $+$ отбор.

Отсюда дешёвый способ собрать предпочтительные пары для псевдо-DPO. DPO обучается на
тройках (промпт, выбранный $y_w$, отвергнутый $y_l$), максимизируя
\[
\log\sigma\!\Big(\beta\log\tfrac{\pi_\theta(y_w|x)}{\pi_{\text{ref}}(y_w|x)}
-\beta\log\tfrac{\pi_\theta(y_l|x)}{\pi_{\text{ref}}(y_l|x)}\Big).
\]
Подавать raw и guided как двух кандидатов на один промпт и ранжировать (частично
автоматически $\varphi$, руками — правки) — чистый способ получить \emph{реально
разных} кандидатов; раз raw иногда лучше, guided иногда — сигнал живой.

\section{Мост между базами: эмбеддинг ConceptNet}

Узлы ATOMIC — события с плейсхолдерами (\textsc{PersonX breaks a cup}), а не
концепты; «кружки» как узла там нет. Мост — извлечь контент-слова события и слинковать
к концепт-узлам ConceptNet. Но \emph{не} bag-of-words/one-hot: они не знают синонимии
(mug $\ne$ cup, разбить $\ne$ сломать), мост выйдет дырявый. Берём \emph{ConceptNet
Numberbatch} — родной эмбеддинг ConceptNet: сначала точный лемма-матч, эмбеддинг-фолбэк.
Тот же эмбеддинг обслуживает $\text{rel}(B\mid q)$ в \eqref{eq:trans} и сидирование
валентности. Одно пространство — три задачи.
